% Route values, deterministic choice, and conditional sorting.
\section{Route Values and Organizational Sorting}
\label{sec:p04-route-sorting}

This section composes the commercial-value kernel with the internal and
entrusted manufacturing technologies defined above. It treats the qualified
manufacturing-service price \(p_m\) as fixed and conjectured. Market
consistency for \(p_m^*\) is imposed in the capacity-market section below.

\subsection{Route values}
\label{subsec:p04-route-values}

For an internally feasible project, the internal route value is
\begin{equation}
  W_i^I(q,m)
  =
  s(q)R\!\left(q,c_I(m,k_i)\right)
  -
  F_I(m,k_i).
  \label{eq:p04-internal-value}
\end{equation}
If \(k_i<\underline{k}(m)\), then \(F_I=+\infty\) and
\(W_i^I=-\infty\), so internal production is unavailable.

At a conjectured CMO capacity price \(p_m\), retained entrusted manufacturing
has value
\begin{equation}
  \begin{aligned}
  W_i^E(q,m;M,p_m)
  ={}&
  s(q)R\!\left(q,c_E(m)\right)
  -F_E(m)-p_m b(m)\\
  &-\mu_E-\tau_E(M).
  \end{aligned}
  \label{eq:p04-entrusted-value}
\end{equation}
The developer remains the authorization holder under \(E\); the qualified
external producer manufactures. Each route cost is subtracted exactly once.
In particular, \(p_m b(m)\) is not contained in \(c_E(m)\).

Let the noncore transfer/out-license outside value be a finite continuous
function \(T(q,m)\), and normalize abandonment or indefinite delay to zero:
\begin{equation}
  W^T(q,m)=T(q,m),
  \qquad
  W^A=0.
  \label{eq:p04-outside-values}
\end{equation}
No separate transfer market is added.

\subsection{Deterministic route choice}
\label{subsec:p04-deterministic-choice}

The optimized value of a planning-stage project is
\begin{equation}
  W_i(q,m;M,p_m)
  =
  \max\left\{
    W_i^I(q,m),
    W_i^E(q,m;M,p_m),
    W^T(q,m),
    W^A
  \right\}.
  \label{eq:p04-optimized-value}
\end{equation}
The route choice is
\begin{equation}
  r_i^*(q,m;M,p_m)
  =
  \arg\max_{r\in\{I,E,T,A\}}
  \begin{cases}
    W_i^I(q,m), & r=I,\\
    W_i^E(q,m;M,p_m), & r=E,\\
    W^T(q,m), & r=T,\\
    W^A, & r=A.
  \end{cases}
  \label{eq:p04-route-choice}
\end{equation}
Continuous project and developer heterogeneity makes exact ties a
measure-zero event, so the choice is single valued almost surely. This is a
deterministic comparison; it generates no probabilistic route share.

When \(M=0\), \(\tau_E(0)=+\infty\) implies \(W_i^E=-\infty\).
At fixed \(p_m\), the binary reform comparison is therefore
\begin{equation}
  \begin{aligned}
  &W_i(q,m;1,p_m)-W_i(q,m;0,p_m)\\
  &\quad=
  \max\!\left\{
    W_i^E(q,m;1,p_m),W_i^I(q,m),W^T(q,m),0
  \right\}\\
  &\qquad\quad-\max\!\left\{W_i^I(q,m),W^T(q,m),0\right\}
  \geq0.
  \end{aligned}
  \label{eq:p04-binary-value-effect}
\end{equation}
The gain is exactly zero, and MAH has zero project-value and route-choice
effect, whenever the newly
feasible entrusted value fails to exceed the best pre-reform route.

\subsection{Internal--entrusted value gap}
\label{subsec:p04-value-gap}

For fixed \(q,m,M,p_m\), define
\begin{align}
  \Delta_{IE}(k_i;q,m,M,p_m)
  &=
  W_i^I(q,m)-W_i^E(q,m;M,p_m) \notag\\
  &=
  s(q)\!\left[
    R\!\left(q,c_I(m,k_i)\right)
    -R\!\left(q,c_E(m)\right)
  \right]
  -F_I(m,k_i) \notag\\
  &\quad
  +F_E(m)+p_m b(m)+\mu_E+\tau_E(M).
  \label{eq:p04-value-gap}
\end{align}
On the internally feasible domain, only the internal technology terms vary
with \(k_i\). Holding \(q,m,M,p_m\) fixed,
\begin{equation}
  \frac{\partial\Delta_{IE}}{\partial k_i}
  =
  s(q)R_c\!\left(q,c_I(m,k_i)\right)c_{I,k}(m,k_i)
  -
  F_{I,k}(m,k_i)
  >0.
  \label{eq:p04-gap-slope}
\end{equation}
The first term is weakly positive because \(s(q)\geq0\),
\(R_c<0\), and \(c_{I,k}<0\); the second is strictly positive because
\(F_{I,k}<0\). Hence the entire derivative is strictly positive, including
the boundary case \(s(q)=0\).

\subsection{Conditional organizational cutoff}
\label{subsec:p04-cutoff}

Consider a finite entrusted barrier and suppose that the continuous value gap
satisfies the endpoint crossing conditions
\[
  \lim_{k\downarrow\underline{k}(m)}\Delta_{IE}(k)<0,
  \qquad
  \lim_{k\to\infty}\Delta_{IE}(k)>0.
\]
Continuity gives existence of a root, and
\eqref{eq:p04-gap-slope} gives strict monotonicity and therefore uniqueness.
Denote the unique root by
\begin{equation}
  \Delta_{IE}\!\left(k^*(q,m;p_m,M);q,m,M,p_m\right)=0.
  \label{eq:p04-cutoff}
\end{equation}
Conditional on both \(I\) and \(E\) dominating transfer and abandonment,
\begin{equation}
  k_i<k^*\Rightarrow r_i^*=E,
  \qquad
  k_i>k^*\Rightarrow r_i^*=I.
  \label{eq:p04-cutoff-sorting}
\end{equation}
If \(T\) or \(A\) dominates, the \(I/E\) cutoff does not determine the
realized route.

For a finite wedge, let the arguments \(\tau_E\) and \(p_m\) vary locally
while holding \(q,m\) fixed. Since
\[
  \frac{\partial\Delta_{IE}}{\partial\tau_E}=1,
  \qquad
  \frac{\partial\Delta_{IE}}{\partial p_m}=b(m)>0,
\]
the implicit-function theorem yields
\begin{equation}
  \frac{\partial k^*}{\partial\tau_E}
  =
  -\frac{1}{\Delta_{IE,k}(k^*)}<0,
  \qquad
  \frac{\partial k^*}{\partial p_m}
  =
  -\frac{b(m)}{\Delta_{IE,k}(k^*)}<0.
  \label{eq:p04-cutoff-derivatives}
\end{equation}
A lower finite institutional barrier raises \(k^*\), expanding the
entrusted-production region \(k_i<k^*\). A higher fixed CMO price lowers
\(k^*\), contracting that region. The second derivative is a fixed-price
comparative static, not an equilibrium derivative through \(p_m^*\).

\subsection{Organizational sorting result}
\label{subsec:p04-sorting-proposition}

Suppose the commercial-value and technology restrictions above hold,
the entrusted barrier is finite, the endpoint crossing conditions hold, and
ties have measure zero. Then the internal-minus-entrusted value gap is strictly
increasing in developer manufacturing capability, a unique finite \(I/E\)
cutoff exists, and the conditional sorting and implicit derivatives in
\eqref{eq:p04-cutoff-sorting}--\eqref{eq:p04-cutoff-derivatives} follow.
The proposition is conditional: it does not displace the valid transfer and
abandonment outside options.

\subsection{Zero-effect and non-tautology boundary}
\label{subsec:p04-zero-effect}

The route comparison is not tautological. Internal value depends on
developer capability through \(c_I\) and \(F_I\); entrusted value instead uses
qualified-external technology, a physical capacity requirement, a retained
holder burden, and the institutional wedge. MAH can nevertheless have zero
effect when \(E\) remains below the pre-reform maximum, when \(T\) or \(A\)
dominates, or when the relevant project--firm pair lies outside the finite
cutoff-crossing region. No smooth random-utility term, inclusive-value object,
or probabilistic route share is used.
